\section{Properties}\label{sec:curve-properties}

The stack proof gives an answer to a gluing question: compatible local
families, and compatible local isomorphisms, come from global ones.
We now ask geometric questions about that stack. These require additional
theorems. Throughout this section $g,n\geq 0$, $2g-2+n>0$, and
$\Mgn$ is the stack over $\CC$ just constructed.

\subsection{What the geometric words ask us to construct}

\begin{Def}[Representability and an atlas]
Start with a morphism $F:\mathcal Y\to\mathcal Z$ of stacks. To test
whether it is \emph{representable by algebraic spaces}, choose any scheme
$S$ and any object $z\in\mathcal Z(S)$, equivalently a map
$S\to\mathcal Z$, and form
\[
\begin{tikzcd}
S\times_{\mathcal Z}\mathcal Y \arrow[r]\arrow[d]
  & \mathcal Y\arrow[d,"F"]\\
S\arrow[r,"z"] & \mathcal Z.
\end{tikzcd}
\]
The requirement is that the fiber product be an algebraic space.
An algebraic space is an \'etale sheaf which has a representable,
surjective \'etale presentation
by a scheme; it allows \'etale coordinate charts even when those charts
do not glue to a scheme by open immersions. In the applications below
these particular fiber products will actually be schemes.

If $F$ is representable, saying that it is smooth, \'etale, or
surjective means that the resulting maps of algebraic spaces to $S$
have that property for every such test. An \emph{atlas} is a
representable surjective map $U\to\mathcal Z$ from a scheme, with the
specified smoothness or \'etaleness condition.
\end{Def}

For example, a map $U\to\Mgn$ is a family of stable pointed curves
over $U$. Given another family $C/S$, a point of
$S\times_{\Mgn}U$ consists of a parameter $u$ in $U$ and an
\emph{isomorphism} from the curve selected by $u$ to the corresponding
fiber of $C/S$. The isomorphism is part of the data. The fiber product
does not merely record that two isomorphism classes agree.

\begin{Def}[Algebraic, Artin, and Deligne--Mumford]
Start with a stack $\mathcal Z$ on the \'etale site of schemes over
$\CC$. In these notes, it is an \emph{algebraic stack}, also called
an \emph{Artin stack}, if its diagonal is representable by algebraic
spaces and it has a smooth atlas $U\to\mathcal Z$ from a scheme.
It is a \emph{Deligne--Mumford stack} if it has an \'etale atlas.
Thus every Deligne--Mumford stack is an algebraic stack in our
convention, since \'etale morphisms are smooth.
\end{Def}

The diagonal has a concrete interpretation here. Pulling
$\Delta:\Mgn\to\Mgn\times\Mgn$ back by two families $C/S,D/S$
gives the sheaf $\Iso_S(C,D)$, with the markings understood.
The earlier sheaf proof showed that this functor is a sheaf.
Representability of the diagonal asks for more: this sheaf must
itself be an algebraic space. An atlas asks for still more: enough
families must be parametrized by schemes in a smooth or \'etale way.
For these conventions and their relation, see Alper
\cite[Definitions 4.1.3--4.1.6 and Theorem 4.2.1, pp.~105--106,
118--119]{Alper}. Some older sources, including Deligne--Mumford,
use ``algebraic stack'' for what we now call a Deligne--Mumford stack.

\subsection{Algebraicity and finite type}

\begin{Prop}\label{prop:curves-algebraic}
The stack $\Mgn$ is algebraic and of finite type over $\CC$.
Its diagonal is affine and of finite type. It has a smooth
surjective atlas $H\to\Mgn$, where $H$ is a scheme of finite
type parametrizing suitable embedded stable pointed curves.
\end{Prop}

\noindent\textbf{Original source and proof strategy.}
Deligne--Mumford \cite[Section 1, pp.~77--78, and Proposition (5.1),
p.~104]{DM} construct the unpointed stack using tricanonical Hilbert
schemes. Knudsen \cite[Theorem 2.7, p.~179]{Knudsen} proves the
pointed theorem over $\operatorname{Spec}\mathbb Z$, using the
unpointed result and induction in the number of markings.
We give a direct Hilbert-atlas argument for pointed stable curves.
Its embedding bound is Alper
\cite[Proposition 6.3.23, p.~231]{Alper}; its smooth-atlas proof is
the embedding-lifting argument of
\cite[Theorem 6.4.6, pp.~238--240]{Alper}, applied here directly
 to our stable category. Although that theorem is written for a
larger category of curves, we do not construct that larger stack.
Our parameter scheme retains some extra choices of embedding;
these choices are allowed in an atlas.

\medskip
\noindent\textbf{Our guided proof.}
The data already available are arbitrary stable pointed families
$(\pi:C\to S,\sigma_1,\ldots,\sigma_n)$. We must construct one
scheme containing enough such families and show that its map to
our stack is representable, smooth, and surjective.

\smallskip
\noindent\emph{Step 1: use a canonical line bundle to bound embeddings.}
Put
\[
 L_C=\omega_{C/S}\Bigl(\sum_{a=1}^n\sigma_a\Bigr).
\]
We use the following \emph{external embedding theorem}: for a stable
pointed family, $L_C^{\otimes3}$ is relatively very ample,
$E_C=\pi_*L_C^{\otimes3}$ is locally free of rank $r=5g-5+3n$,
and formation of this vector bundle commutes with arbitrary base
change. These conclusions follow from the fiberwise vanishing and
very ampleness theorem together with cohomology and base change;
see Alper \cite[Exercise 6.3.18 and Proposition 6.3.23,
pp.~230--231]{Alper}. The fiberwise unpointed theorem is proved in
Deligne--Mumford \cite[Theorem (1.2) and its corollary,
pp.~77--78]{DM}. We take the pointed very ampleness result as an
input rather than redo its separation-of-points-and-tangents proof.

The rank and Hilbert polynomial have a useful numerical explanation.
On any geometric fiber, let $\ell=2g-2+n>0$. Then
$\deg L_C=\ell$ and $\deg L_C^{\otimes3}=3\ell$. Vanishing of
$H^1(C,L_C^{\otimes3})$ and Riemann--Roch give
\[
 r=3\ell+1-g=5g-5+3n,
 \qquad P(m)=\chi(C,L_C^{\otimes3m})=3\ell m+1-g.
\]
These numbers depend only on $g,n$, not on the particular curve.
The evaluation of sections gives a closed immersion
$C\hookrightarrow\mathbb P_S(E_C)$, where our convention for
$\mathbb P(E_C)$ parametrizes one-dimensional quotients of $E_C$.
Locally on $S$, choosing a basis of $E_C$ identifies the ambient
bundle with $\mathbb P^{r-1}_S$.

\smallskip
\noindent\emph{Step 2: represent the diagonal.}
Take two stable pointed families $C/S,D/S$. After trivializing
$E_C,E_D$ locally over $S$, their canonical embeddings place them
in the same projective space. An isomorphism $C\to D$ induces an
isomorphism of their canonical line bundles, hence a unique
projective linear transformation of these complete embeddings.
Conversely, a projective transformation carrying the embedded
curve and every marking to the corresponding data restricts to
an isomorphism of pointed families.

Thus $\Iso_S(C,D)$ is locally represented by a closed subscheme
of $\PGL_{r,S}$: the transformation must send one Hilbert point
 to the other, and must send each marked point to its partner.
Equality of Hilbert points is closed because the Hilbert scheme
is separated; the point conditions are closed because projective
space is separated. These are finitely presented conditions, since
the Hilbert scheme and projective space are of finite presentation
over $\CC$. The local
representing schemes glue by descent for affine schemes, so
$\Iso_S(C,D)$ is an affine scheme of finite presentation over $S$.
This proves the claimed representability and finiteness of the
diagonal. It also implies that a map from any scheme to our
stack is representable: its test fiber products are obtained by
base changing that diagonal. Notice how much stronger this is
than merely saying that $\Iso_S(C,D)$ is a sheaf.

\smallskip
\noindent\emph{Step 3: choose a finite type parameter scheme.}
The \emph{Hilbert scheme theorem} is an external input:
closed subschemes of $\mathbb P^{r-1}_S$ flat and finitely presented
over $S$, with fiber Hilbert polynomial $P$, are represented by a
projective scheme $\operatorname{Hilb}^{P}(\mathbb P^{r-1}_{\CC})$.
See Alper \cite[Theorem 2.1.2, p.~48]{Alper}.
Adding $n$ points and requiring each to lie on the universal curve
gives a closed incidence scheme in
\[
 \operatorname{Hilb}^{P}(\mathbb P^{r-1}_{\CC})
                      \times(\mathbb P^{r-1}_{\CC})^n.
\]
Inside it, take the open locus $H$ where the pointed fibers are
stable and $H^1(C_s,\mathcal O_{C_s}(1))=0$.
Openness of stability is Alper
\cite[Proposition 6.3.24, p.~232]{Alper}; the vanishing locus is
open by upper semicontinuity of cohomology. Both are substantial
standard inputs about projective flat families. The Hilbert
polynomial already fixes the arithmetic genus to be $g$.
Since the incidence scheme is of finite type over the noetherian
field $\CC$, its open subscheme $H$ is of finite type.

The universal family defines $a:H\to\Mgn$ by forgetting the
embedding. It is representable by Step 2. It is surjective because
any stable family is, locally on its base, embedded by $L_C^3$
as in Step 1, and these embeddings satisfy the defining conditions
of $H$. We deliberately allow other embeddings in $H$ as well;
being an atlas does not require a unique embedding for each curve.

\smallskip
\noindent\emph{Step 4: why forgetting the embedding is smooth.}
Smoothness says that an embedding can be chosen compatibly when a
family undergoes a small infinitesimal variation. We apply the
\emph{infinitesimal criterion for smoothness} to the representable
map $a$, following Alper
\cite[Theorem 4.7.1, pp.~144--145; proof of Theorem 6.4.6,
pp.~238--240]{Alper}. The finite-presentation hypothesis of the
criterion also holds. For a test family $C/S$, the scheme
$S\times_{\Mgn}H$ is the isomorphism scheme between the two
families over $S\times H$: the pullback of $C/S$ and the pullback
of the universal family over $H$. Step 2 makes this scheme
finitely presented over $S\times H$, and $S\times H$ is finitely
presented over $S$. Thus $a$ is representable and locally of finite
presentation before we use the lifting calculation.

Let $A'\twoheadrightarrow A$ be a small extension of local Artinian
$\CC$-algebras, with residue field $K$ and kernel $J$ killed by the
maximal ideal of $A'$. We are given a stable pointed family $C'/A'$
and an embedding of its restriction $C=C'\times_{A'}A$ in
$\mathbb P^{r-1}_A$, representing a point of $H(A)$. We must lift
that embedding to $C'$:
\[
\begin{tikzcd}
 C\arrow[r,hook]\arrow[d,hook] & \mathbb P^{r-1}_A\arrow[d,hook]\\
 C'\arrow[r,dashed,hook] & \mathbb P^{r-1}_{A'}.
\end{tikzcd}
\]
The embedding on $C$ is given by $M=\mathcal O_C(1)$ and its
$r$ coordinate sections.

A line bundle on $C$ lifts to $C'$ because its lifting obstruction
lies in $H^2(C_K,\mathcal O_{C_K}\otimes_KJ)=0$.
This elementary line-bundle lifting theorem follows from
$1+\mathcal I\to\mathcal O_{C'}^*\to\mathcal O_C^*$ for the
square-zero ideal $\mathcal I$; see Alper
\cite[Proposition C.2.11 and Remark C.2.12, p.~602]{Alper}.
Choose a lift $M'$. The restriction sequence is
\[
 0\longrightarrow M_K\otimes_KJ
   \longrightarrow M'\longrightarrow M\longrightarrow0.
\]
The defining vanishing $H^1(C_K,M_K)=0$ implies that every
coordinate section of $M$ lifts to $M'$. The lifted sections
still generate $M'$ by Nakayama's lemma. They define a morphism
$C'\to\mathbb P^{r-1}_{A'}$ lifting the given embedding.

It is a closed immersion. Indeed it is proper, its only geometric
fiber is a closed immersion and hence it is quasi-finite, so it
is finite. The cokernel of the map from the ambient structure
sheaf to its finite pushforward restricts to zero over $A$ and
therefore vanishes by Nakayama's lemma. The resulting embedded
family is flat over $A'$ because it is the given family $C'$.
Its marked sections were already given and acquire their images
under this embedding. Its special fiber lies in $H$, so the
lift also lies in the open locus $H$. This completes the required
lifting. The infinitesimal criterion now proves that $a$ is smooth.

We have constructed a representable smooth surjective map from
the scheme $H$ to $\Mgn$, and have represented the diagonal.
These are exactly the requirements for algebraicity. Finally,
finite type of $H$ proves that $\Mgn$ is quasi-compact and
locally of finite type over $\CC$, which is the asserted finite
type property.

\begin{Ex}[Genus two and extra embedding choices]
For $(g,n)=(2,0)$ we have $\ell=2$, $r=5$, and $P(m)=6m-1$.
Every stable genus-two curve has a tricanonical embedding in
$\mathbb P^4$. Thus the Hilbert scheme for that one polynomial
already contains embeddings of every stable genus-two curve.
The atlas retains choices of an embedding and coordinates;
$\overline{\mathcal M}_2$ retains the abstract curve and its
isomorphisms. Extra parameters in an atlas are harmless: a
smooth atlas need not have relative dimension zero.
\end{Ex}

\subsection{Why the diagonal is unramified}

\begin{Prop}\label{prop:curves-dm}
Under the same hypotheses, $\Mgn$ is Deligne--Mumford.
Every geometric stable pointed curve has a finite reduced
automorphism group scheme.
\end{Prop}

\noindent\textbf{Original source and proof strategy.}
Deligne--Mumford \cite[Lemma (1.4), pp.~80--81, Theorem (1.11),
p.~84, and Proposition (5.1), p.~104]{DM} eliminate infinitesimal
automorphisms and construct the unpointed Deligne--Mumford stack.
Knudsen's pointed theorem is \cite[Theorem 2.7, p.~179]{Knudsen}.
The pointed calculation is explained by Alper
\cite[Proposition 6.3.26, pp.~232--235, and Theorem 6.4.16,
p.~241]{Alper}. The key geometric fact is that every component of
the normalization has enough special points to eliminate regular
vector fields fixing those points.

\medskip
\noindent\textbf{Our guided proof.}
We already have algebraicity. The additional construction we need
is an \'etale atlas. We use the following \emph{Deligne--Mumford
criterion} as an external theorem: an algebraic stack has an
\'etale atlas if and only if its diagonal is unramified, equivalently
if every geometric stabilizer is discrete and reduced. In our
finite type situation these stabilizers will be finite.
See Alper \cite[Theorem 4.6.4, p.~143]{Alper}.

Fix an algebraically closed extension $K/\CC$ and a stable pointed
curve $(C,p_1,\ldots,p_n)$ over $K$. An \emph{infinitesimal
automorphism} starts with the constant family over
$K[\epsilon]/(\epsilon^2)$ and is an automorphism that restricts
to the identity when $\epsilon=0$. Locally on functions it has
the form $a\mapsto a+\epsilon D(a)$, where the multiplication
rule forces $D$ to be a derivation. Fixing a marked point forces
the vector field $D$ to vanish there. Thus the tangent space of
the automorphism scheme at its identity is
\[
 H^0\bigl(C,\mathcal Hom(\Omega_C,\mathcal O_C(-\sum p_a))\bigr).
\]

Let $\nu:\widetilde C\to C$ be the normalization. On a nodal
curve a regular vector field lifts to the normalization and
vanishes at each branch over every node. This can be checked
locally in $K[[x,y]]/(xy)$ by applying a derivation to $xy=0$;
see Alper \cite[Exercise 6.2.22, p.~221]{Alper}.
Write $C_v$ for the smooth connected components of $\widetilde C$,
$h_v$ for their genera, and $B_v$ for the marked points and
branches of nodes on $C_v$. If $m_v=\deg B_v$, stability says
\[
                  2h_v-2+m_v>0.
\]
Consequently
\[
 \deg T_{C_v}(-B_v)=2-2h_v-m_v<0.
\]
A line bundle of negative degree on a smooth projective connected
curve has no nonzero section. Every component vector field
therefore vanishes, so the infinitesimal automorphism vanishes.

The automorphism scheme is an affine group scheme of finite type
by the preceding diagonal calculation. Its zero tangent space at
the identity implies that the identity is an isolated reduced
point: in its noetherian local ring the equality
$\mathfrak m/\mathfrak m^2=0$ implies $\mathfrak m=0$ by
Nakayama's lemma. Translation gives the same conclusion at every
geometric point. A zero-dimensional reduced scheme of finite
type over $K$ has finitely many points. Thus the stabilizer is
finite and reduced, and the Deligne--Mumford criterion applies.

In ordinary language, there can be discrete symmetries of the
curve, but there is no continuously varying infinitesimal symmetry.
For a genus-two hyperelliptic curve, the involution remains an
automorphism; it does not prevent the stack from being
Deligne--Mumford. It does prevent us from identifying that moduli
problem with a scheme which has the same fiber groupoids.

\begin{Ej}\label{ex:curve-dm}
Check the three cases $h_v=0$, $h_v=1$, and $h_v\geq2$ in the
inequality above. Why would a smooth unmarked elliptic curve fail
this argument? Explain why ``finite automorphism groups'' and
``trivial automorphism groups'' are different requirements.
\end{Ej}

\subsection{Properness and the meaning of a limit}

\begin{Prop}\label{prop:curves-proper}
The Deligne--Mumford stack $\Mgn$ is separated and proper over
$\CC$.
\end{Prop}

\noindent\textbf{Original source and proof strategy.}
Deligne--Mumford \cite[Section 2 and Theorem (5.2),
pp.~104--105]{DM} prove the unpointed properness theorem using
stable reduction and uniqueness of stable limits.
Knudsen \cite[Theorem 2.7, p.~179]{Knudsen} proves the pointed
version. A later detailed explanation in characteristic zero is
Alper \cite[Theorem 6.5.1 and Section 6.5.1, pp.~242--247;
Proposition 6.5.20 and Theorem 6.5.23, pp.~251--252]{Alper}.
Alper's written uniqueness proof is unpointed; the marked
extension is assigned in Exercise 6.5.22. We take the full pointed
existence and uniqueness statements from the established pointed
theorem as external inputs. We do not claim to have supplied a
new proof of stable reduction here.

\medskip
\noindent\textbf{Our guided proof.}
For an algebraic stack, \emph{separated} means that its diagonal
is proper. In terms of curves this is a statement about extending
\emph{isomorphisms} between families. \emph{Proper over $\CC$}
means separated, of finite type, and universally closed over
$\CC$. The last condition is a completeness condition: limits
cannot escape after any change of base. To test it effectively we
use the stack valuative criterion
\cite[Theorem 4.8.7, pp.~149--150]{Alper}.

Here is the data in that criterion. Let $R$ be a discrete valuation
ring over $\CC$, let $K$ be its fraction field, and let
$(C_K,p_1,\ldots,p_n)$ be a stable pointed curve over $K$.
Think of $\operatorname{Spec}K$ as the punctured base and of
the closed point of $\operatorname{Spec}R$ as the missing limit.
The \emph{pointed stable reduction theorem} supplies, after a
finite extension of the fraction field and a corresponding
extension of discrete valuation rings $R\subset R'$, a stable
pointed family $C'\to\operatorname{Spec}R'$ whose generic fiber
is the prescribed base change of $C_K$.

The second external input is \emph{uniqueness of stable limits}:
if $C/R$ and $D/R$ are two stable pointed families, then every
specified pointed isomorphism $C_K\to D_K$ extends to a unique
pointed isomorphism $C\to D$. The word ``specified'' matters.
An object may have automorphisms; uniqueness concerns extension
of one given generic isomorphism, not uniqueness of all possible
isomorphisms between two limits.

Apply uniqueness to the valuative criterion for the representable
diagonal. The diagonal is already affine and of finite type, hence
quasi-compact and separated. The criterion proves that it is
proper, so $\Mgn$ is separated. In particular its diagonal is
finite, since an affine proper morphism is finite.
Apply existence to the valuative criterion for the finite type
stack $\Mgn\to\operatorname{Spec}\CC$. The permitted finite
base extension in this criterion is exactly what stable reduction
provides. The criterion proves universal closedness and hence
properness. The hypotheses of the criterion have all been supplied:
algebraicity, finite type, and quasi-separatedness came from the
Hilbert presentation, and separatedness was just checked.
Local noetherianness follows from finite type over $\CC$.

Why does stable reduction involve geometry beyond descent?
An \'etale cover sees every point of the base. The inclusion
$\operatorname{Spec}K\hookrightarrow\operatorname{Spec}R$
misses the closed point, so the stack condition cannot manufacture
a curve over that point. Stable reduction is the theorem that
does so, using the particular geometry of nodal stable curves.

\begin{Ex}[A genus-two stable boundary point]
Let $E_1,E_2$ be smooth elliptic curves and identify one point
on each to form $C=E_1\cup_qE_2$. Its arithmetic genus is two,
and each elliptic component has one node as a special point,
so $C$ is stable. A smoothing of this node has local equation
$xy=t$; the smooth fibers have genus two. The stable fiber at
$t=0$ belongs to $\overline{\mathcal M}_2$, while it does not
belong to the open smooth-curve stack $\mathcal M_2$.
The deformation theorem below supplies such smoothings. Uniqueness
of stable limits shows that this punctured family cannot instead
acquire a smooth special fiber after finite base extension.
This illustrates why $\mathcal M_2$ is not proper and why the
bar in $\overline{\mathcal M}_2$ matters.
\end{Ex}

\subsection{Smoothness and actual dimension}

\begin{Prop}\label{prop:curves-smooth}
The stack $\Mgn$ is smooth over $\CC$ and has pure dimension
$3g-3+n$.
\end{Prop}

\noindent\textbf{Original source and proof strategy.}
Deligne--Mumford \cite[Lemmas (1.3)--(1.4), Proposition (1.5),
pp.~79--82, and Theorem (5.2), pp.~104--105]{DM} analyze
deformations of nodal curves and obtain the unpointed smoothness
statement. Knudsen \cite[Theorem 2.7, p.~179]{Knudsen} gives the
pointed theorem. Alper \cite[Propositions 6.3.26 and 6.3.28,
pp.~232--235; Theorem 6.4.16, p.~241]{Alper} computes the
pointed deformation spaces and their dimension in detail.
We explain the count on the normalization and state the deformation
theorem used to pass from that count to smoothness.

\medskip
\noindent\textbf{Our guided proof.}
Smoothness of a Deligne--Mumford stack over $\CC$ can be tested
on an \'etale atlas. Its dimension is the dimension of that atlas
locally, because an \'etale map has relative dimension zero.
To prove smoothness we use the infinitesimal lifting criterion:
a small deformation over an Artinian base must lift when the base
is enlarged by a square-zero ideal. A deformation of a pointed
curve means a family together with an identification of its special
fiber with that fixed pointed curve.

Our \emph{external nodal deformation theorem} has two parts.
First, deformations of a proper nodal curve with distinct smooth
marked points have no obstruction to such infinitesimal lifting.
Second, the vector space of first-order deformations fits in
\begin{equation}\label{eq:curve-deformations}
 0\longrightarrow
 \bigoplus_v H^1(C_v,T_{C_v}(-B_v))
 \longrightarrow \operatorname{Def}(C,p_1,\ldots,p_n)
 \longrightarrow \bigoplus_{q\text{ a node}} T_q
 \longrightarrow0,
\end{equation}
where each $T_q$ is a one-dimensional smoothing space.
Choosing local coordinates at the two branches identifies it
with $K$, with deformation equation $xy=\epsilon a$.
That identification depends on coordinates, but its dimension
does not. The theorem is Alper's
\cite[Propositions 6.3.26 and 6.3.28, pp.~232--235]{Alper}.
The local calculation behind unobstructedness is that a node is
a single equation $xy=0$ whose deformation parameter can be
lifted, while a curve has no degree-two coherent cohomology to
obstruct the global gluing. The theorem makes that local-to-global
argument precise; we are not developing a general obstruction
theory in these notes.

The left side of \eqref{eq:curve-deformations} varies the smooth
normalization components, with every marked point and every
branch of a node remembered. Such variations keep the nodes
unsmoothed. The right side allows one smoothing parameter for
each node. Surjectivity says that these local smoothing parameters
can be realized by a global first-order deformation.

Let $V$ be the number of normalization components and let
$\delta$ be the number of nodes. Riemann--Roch and the vanishing
of $H^0(C_v,T_{C_v}(-B_v))$ proved above give
\[
 h^1(C_v,T_{C_v}(-B_v))=3h_v-3+m_v.
\]
There are $n$ marked points and two branches per node, so
$\sum_vm_v=n+2\delta$. The normalization sequence for functions
gives the genus formula
\[
                  g=\sum_vh_v+\delta-V+1.
\]
Taking dimensions in \eqref{eq:curve-deformations} now gives
\begin{align*}
 \dim\operatorname{Def}(C,p_1,\ldots,p_n)
 &=\sum_v(3h_v-3+m_v)+\delta\\
 &=3\sum_vh_v-3V+n+3\delta\\
 &=3g-3+n.
\end{align*}
There are no infinitesimal automorphisms, by
Proposition~\ref{prop:curves-dm}, so this is the tangent dimension
on an \'etale chart, without any extra infinitesimal symmetry
directions to remove. Unobstructedness and the infinitesimal
lifting criterion show that the charts are smooth; the constant
tangent dimension then proves the dimension assertion.

\begin{Ex}[Checking the dimension at a nodal genus-two curve]
For $C=E_1\cup_qE_2$, each normalization component is elliptic
with one special point. It contributes
$3\cdot1-3+1=1$ parameter. The node contributes one more.
The total is $1+1+1=3$, just as at a smooth genus-two curve.
The locus where this particular node remains unsmoothed has
only the first two parameters. A singular member of the family
does not force the moduli stack itself to be singular.
\end{Ex}

\begin{Ej}\label{ex:curve-dimension}
Glue two copies of $\mathbb P^1$ along three distinct pairs of
points, with no markings. Check stability and genus. Use
\eqref{eq:curve-deformations} to count the dimension of its
deformation space. Which parameters come from the components,
and which come from the nodes?
\end{Ej}

\subsection{What depends on the hypotheses}

All four propositions hold for every stable pair $g,n$ in the
setup of this chapter. The historical theorems actually prove
the corresponding proper smooth Deligne--Mumford statements over
$\operatorname{Spec}\mathbb Z$; we used their base change to
$\CC$, with the characteristic-zero expositions indicated.

If $2g-2+n\leq0$, the stable pointed moduli problem as defined
here has no geometric objects. One can still speak of its empty
stack, but not of a nonempty smooth stack of dimension $3g-3+n$.
If one instead drops stability while keeping those curves, the
conclusions change: an unmarked $\mathbb P^1$ has the
positive-dimensional group $\PGL_2$ of automorphisms, and an
unmarked elliptic curve has translations. Thus the argument for
the Deligne--Mumford property would fail.

Finally, smoothness belongs to the \emph{stack}; it is not an
automatic assertion about its coarse moduli space. Finite group
quotients of smooth charts can be singular. We have not needed
the existence or projectivity of a coarse moduli space for either
the stack proof or the properties proved here.
